\section{Background}
\label{sec:background}

\subsection{Agent Lifecycle Management in Multi-Agent Systems}
\label{sec:agent-lifecycle}

The lifecycle of an autonomous agent spans four canonical phases: \textbf{genesis} (creation and identity issuance), \textbf{execution} (task processing and tool invocation), \textbf{spawning} (delegation to child agents), and \textbf{decommission} (graceful termination and state archival). When agents operate in systems of two or more, the management of these phases becomes non-trivial: coordination protocols must enforce that spawned agents inherit scoped authority rather than unbounded privilege, that execution remains within the mandate granted by the parent, and that termination closes all outstanding obligations without leaving orphaned processes.

Multi-agent systems (MAS) have been studied since the mid-1990s as a paradigm for distributed problem-solving \cite{wooldridge1995,stone2000}. Early agent frameworks treated delegation as a messaging contract between peer processes, with no architectural constraint on chain depth or scope inheritance \cite{ferber1996}. This peer-model proved difficult to scale because accountability was diffuse: if a sub-agent acted outside its mandate, the parent could claim ignorance and the system had no forensic mechanism to establish otherwise. 

Modern LLM-based multi-agent systems have revived these concerns with new urgency. Systems such as AutoGen \cite{autogen2023}, LangChain \cite{langchain2023}, and CrewAI \cite{crewai2024} expose rich APIs for agent spawning but treat the delegation relationship as a software configuration rather than an accountability primitive. The consequence is that spawned agents in these systems carry no cryptographically verifiable record of their parentage, scope, or authorized depth. Drift is structurally possible because the chain of delegation is implemented as a mutable runtime annotation, not as an immutable architectural constraint.

Recent proposals address this gap directly. The Agent Lifecycle Protocol (ALP) defines six canonical lifecycle events---Genesis, Fork, Migration, Retraining, Succession, and Decommission---along with formal state machine semantics that govern transitions between them \cite{agent-lifecycle-protocol}. ALP introduces a genealogical registry that tracks both genetic lineage (model and architecture provenance) and epigenetic lineage (configuration, memory, and reputation history), enabling forensic reconstruction of the delegation tree after an incident. The Agent Identity, Trust and Lifecycle Protocol (AITLP) extends this with formal mandate enforcement, drift detection, and revocation mechanisms that operate at the level of the operating system layer rather than the application layer \cite{draft-larsson-aitlp}. Together, these efforts establish that agent lifecycle management is not a deployment detail but an accountability architecture: the design decisions made at provisioning time determine what is and is not auditable at runtime.

\subsection{Accountability and Provenance in Distributed Agent Systems}
\label{sec:accountability-provenance}

Accountability in distributed systems requires that every action can be attributed to the principal that authorized it, and that the chain of authorization is tamper-evident and non-repudiable. For single-agent systems this is straightforward: the agent is the principal. For multi-agent systems the challenge compounds with chain depth: each spawn event introduces a new principal, and the question of whether the child is acting within the parent's scope must be answerable by reference to a verifiable record, not by inference from runtime behavior.

The W3C PROV standard provides a foundational framework for provenance tracking in distributed workflows \cite{provw3c}. PROV defines entities, activities, and agents with attributed generation, usage, and communication relationships that enable post-hoc reconstruction of causal chains. Extending PROV to agentic workflows, PROV-AGENT adapts its model to capture agent-centric metadata---prompts, responses, decisions, and tool invocations---and links them with workflow-level provenance using the Model Context Protocol (MCP) \cite{prov-agent}. PROV-AGENT enables near real-time provenance capture and querying across federated, edge, cloud, and HPC environments, demonstrating that end-to-end accountability is achievable in heterogeneous agent deployments.

Cryptographic chain-of-custody protocols raise the bar further by binding human authorization to agent actions through an append-only, per-hop signed record. The Human Delegation Provenance (HDP) protocol introduces a token-based mechanism where the human principal's authorization, scope, and session binding are captured at issuance; each agent hop adds a signed cryptographic record to an append-only chain \cite{hdp2026}. Verification is offline-capable, requiring only the issuer's public key and session identifier. HDP addresses a critical accountability gap in existing agentic deployments where task descriptions and delegations lack cryptographic binding to the original human authorization.

The Agent Passport System (APS) models authority as a lattice element and delegation as a monotone function on that lattice, with seven dimensions of scope attenuation (scope, spend, depth, time, reputation, values, reversibility) enforced strictly so that delegated authority can narrow but never expand \cite{aps2025}. Cascade revocation in APS invalidates all descendants when a parent certificate is revoked, providing strong provenance control. The Warrant Certificate Authorities (WCA) framework complements this with auditable data provenance for tool-call chains, establishing that each tool invocation can be traced back to the warrant that authorized it \cite{wca2025}. The Intent Provenance Protocol (IPP) further extends this model to capture the semantic intent behind agent actions, not just their mechanical execution \cite{ipp2026}. Together, these protocols establish that accountability requires more than logging: it requires a cryptographically enforced chain of custody from human principal to terminal action.

PEDIGREE extends SPIFFE-style workload identity to multi-agent systems with per-hop cryptographic chaining (\texttt{parent\_chain}) that ensures traceable, non-repudiable lineage of sub-agents \cite{pedigree}. PEDIGREE combines operator-controlled ceilings with per-parent mandate narrowing ($\text{ceiling} \cap \text{mandate}$) to prevent escalation, and enforces scope attenuation at both mint and verify time. This dual enforcement model is architecturally closer to the BX3 approach than any prior work: it recognizes that accountability is a structural property enforced by the conjunction of multiple independent constraint layers, not a property that can be retrofitted through logging alone.

\subsection{Fixed vs. Mutable Delegation Chains}
\label{sec:fixed-vs-mutable}

The distinction between fixed and mutable delegation chains is central to the problem of agentic accountability. A \textbf{fixed delegation chain} is established at agent provisioning and is structurally immutable for the agent's lifetime: the parent reference, authorized scope, and depth bound are sealed into the agent's identity record and cannot be altered without creating a new agent identity. A \textbf{mutable delegation chain} permits these parameters to be updated at runtime, either by the parent agent, by an administrator, or by the agent itself. Mutable chains offer greater flexibility but create a structural vulnerability: an agent that can modify its own delegation parameters can also attenuate the evidence of its own authority, making post-hoc accountability unreliable.

Mutable delegation is the norm in contemporary LLM-based agent frameworks. In AutoGen, LangChain, and CrewAI, the parent-child relationship is a software configuration that can be changed through API calls at runtime. This is architecturally equivalent to giving the agent administrative access to its own authorization record. When something goes wrong---a drifted agent, a prompt injection, an unauthorized tool call---the system cannot establish with certainty what the agent's authorized scope was at the time of the action, because that scope was mutable and the record may have been updated after the fact.

Fixed delegation chains address this at the architecture level. The key invariant is \textbf{immutability of the parent pointer}: once an agent is provisioned, its \texttt{parent} field is sealed by theFact Layer and cannot be modified for the agent's operational lifetime. The root human sets this field to $\bot$ (null); every other agent receives it at provisioning from its parent agent. The invariant guarantees that the delegation chain from any agent to its human root is deterministically recoverable at any point in time, and that no agent can obscure its lineage by retroactive modification of its parent reference. This is not a policy control or an access restriction; it is a structural property enforced by the cryptographic identity layer. The fixed chain model was anticipated in early multi-agent literature \cite{padgham2001} but was not implementable without a practical cryptographic identity substrate for agents. The emergence of agent identity standards (SPIFFE SVIDs, DIDs, Ed25519 agent passports) now makes it feasible \cite{aps2025,draft-beyer-agent-identity-architecture-00}.

Dynamic (mutable) delegation chains have legitimate use cases, particularly in adaptive orchestration systems where task requirements change during execution. ROMA introduces a hierarchical, model-agnostic decomposition framework where agents adopt modular roles (Atomizer, Planner, Executor, Aggregator) and the delegation topology can evolve as sub-task dependencies are resolved \cite{roma2025}. AgentSpawn demonstrates runtime spawning of specialized child agents guided by real-time task complexity analysis, with adaptive delegation depth \cite{agentspawn2026}. FIRMHIVE treats delegation as an executable primitive encoded as actionable steps rather than static handoffs, enabling per-agent mutable chains that remain traceable \cite{firmhive2025}. TDAG similarly decomposes tasks dynamically and generates reconfigurable subagents, proactively mitigating error propagation through structured delegation \cite{tdag2025}. These systems demonstrate that mutable delegation and traceable provenance are not inherently incompatible: mutable delegation becomes accountable when every modification to the delegation chain is itself recorded as a signed event in an immutable forensic ledger. The distinction that matters is not fixed vs. mutable \emph{per se} but whether the mutation events are cryptographically sealed or left as unguarded runtime state.

The Agent Execution Records (AER) framework formalizes this by treating provenance not as metadata but as a first-class execution primitive, coequal with prompts and tool calls \cite{aer2025}. AER records are sealed at the time of action and form an immutable audit trail. The critical insight is that fixing the delegation chain at provisioning while permitting mutable execution within that scope produces a delegation model that is simultaneously adaptable and accountable. The scope can evolve; the chain cannot.

\subsection{Hierarchical Task Decomposition in AI Agents}
\label{sec:htd}

Hierarchical task decomposition (HTD) is the practice of decomposing a complex, high-level goal into a tree or directed acyclic graph (DAG) of subtasks, each assigned to a specialized agent or agent team. HTD is foundational to scalable multi-agent systems because it enables parallel work distribution, domain specialization, and result aggregation across long-horizon tasks. The canonical HTD model in classical AI is Hierarchical Task Network (HTN) planning \cite{htn1993,htn2000}, where domain methods decompose compound tasks into primitive actions with enforced ordering constraints. Modern LLM-agent systems have adopted and extended this model, replacing hand-coded domain methods with learned decomposition driven by Large Language Model reasoning \cite{hiplan2025,reactree2025,structuredagent2025}.

The challenge in LLM-based HTD is that decomposition is non-deterministic: the same goal can be decomposed in multiple functionally distinct ways, and the decomposition strategy determines downstream accountability. If a high-level goal is decomposed into subtasks assigned to child agents, the parent must retain enough context to authorize each subtask scope, monitor each subtask's execution, and aggregate results. If any subtask fails or drifts, the parent must detect the failure, decide whether to respawn or reassign, and record the recovery action. Systems that perform HTD without structuring the parent-child relationship as an immutable accountability primitive risk the same drift and orphaning problems described in Section \ref{sec:fixed-vs-mutable}.

DeepAgent demonstrates HTD using a Hierarchical Task DAG (HTDAG) that recursively decomposes goals while encoding dependencies and constraints in a dynamic graph structure \cite{deepagent2025}. The system operates in a continuous planner-executor loop that adapts both the decomposition and the execution as circumstances evolve. AgentOrchestra implements a central planning agent that decomposes objectives and delegates sub-tasks to specialized sub-agents with MCP-driven tool management, enabling dynamic tool creation and reuse across the hierarchy \cite{agentorchestra2025}. ReCAP introduces recursive context-aware reasoning where each level of decomposition carries its own epistemic context, enabling the system to reason about its own reasoning \cite{recap2025}. These systems demonstrate that HTD is both necessary for scalable task solving and a natural point of integration for accountability enforcement: the decomposition structure is the delegation structure, and the delegation structure is the accountability structure.

The critical design question is where to enforce the accountability constraint relative to the decomposition. Bottom-up approaches (enforce at execution) require runtime monitoring that is expensive and can be circumvented by compromised agents. Top-down approaches (enforce at decomposition) are more robust because they reject invalid spawns before they activate, using theFact Layer pre-commit hook to validate the chain against the Bounds Engine constraints. The BX3 Framework adopts a top-down approach: every spawn event is validated against three independent layers before the child agent enters operational state.

\subsection{The BX3 Framework as Substrate for Immutable Parent Chains}
\label{sec:bx3-framework}

The BX3 (Behavior, Execution, and eXchange) Framework is a universal architecture for accountable autonomous systems that provides three independent functional layers---Purpose Layer, Bounds Engine, and Fact Layer---operating in concert to enforce that every agent in any spawn chain is, by architectural construction, traceable to a human principal, bounded by a finite delegation depth, and recorded in an immutable forensic ledger before entering operational state \cite{bx3framework,bx3framework2026}. BX3 is not a protocol or a policy convention; it is an architectural specification that defines how the three layers interact at every spawn event, every execution gate, and every state transition.

The \textbf{Purpose Layer} encodes the terminal purpose and operational scope of every agent. It serves as the human-anchored accountability anchor: every machine agent has a Purpose Anchor that terminates at a human principal who is accountable for the agent's actions. The Purpose Layer participates actively in the runtime enforcement checkpoint (Truth Gate) by providing the normative reference against which all agent behavior is verified. Before any action is sealed to theFact Layer, the Truth Gate confirms that the proposed action is consistent with the agent's declared Purpose Anchor. An action that cannot be reconciled with the Purpose Anchor fails the Truth Gate and is rejected, preventing drift at the execution level \cite{bx3framework}.

The \textbf{Bounds Engine} enforces structural constraints on every spawn event before activation occurs. It evaluates spawn requests against four bounds categories: delegation depth ($D_{\max}$), scope width ($S_{\max}$), authority bound, and temporal bound ($T_{\text{bounds}}$). The Bounds Engine is an independent, deterministic constraint checker that is structurally isolated from the Purpose Layer: the Purpose Layer defines \emph{what} an agent is authorized to do; the Bounds Engine defines \emph{how many} levels deep, \emph{how many} children, and \emph{for how long} the delegation chain may extend. No spawn event can activate without passing both the Purpose Layer's intent check and the Bounds Engine's structural check. This separation of concerns is architecturally essential: a compromised Purpose Layer cannot bypass depth limits, and a compromised Bounds Engine cannot authorize out-of-scope actions \cite{bx3framework2026}.

The \textbf{Fact Layer} is the deterministic execution firewall: an append-only, cryptographically sealed audit log that records every spawn event, purpose clause binding, Truth Gate passage or rejection, and agent state transition. Every entry is linked to its predecessor via SHA-256 hash chaining, forming an unbroken cryptographic record analogous to a blockchain's transaction ledger \cite{nakamoto2008bitcoin}. The Fact Layer performs three forensic functions: (1) it provides structural immutability---no entry can be modified, deleted, or inserted after the fact without breaking the hash chain; (2) it provides forensic completeness---every spawn event and every Truth Gate transition is recorded before the next event can occur; and (3) it provides non-repudiation---the chain serves as legal evidence that a specific agent acted within a specific scope at a specific time \cite{bailoutprotocol2026}.

The integration of these three layers produces an immutable parent chain architecture that differs fundamentally from logging-based accountability approaches. In BX3, the parent pointer is not a mutable metadata field; it is sealed into the agent's cryptographic identity record at the moment of provisioning and remains structurally inseparable from that identity for the agent's entire operational lifetime \cite{sovereignagent}. The root human sets the null-parent case ($\parent(h) = \bot$); every machine agent receives its parent pointer from its parent agent at provisioning. Any spawn event that would produce a machine agent with a null parent is rejected by theFact Layer pre-commit hook, enforcing the human-grounding invariant: every machine agent in a BX3-compliant system has a traceable lineage to a human principal \cite{bx3framework2026}.

The Recursive Spawning Protocol (RSP), the central topic of this paper, is the operationalization of BX3 Pillar 2 (Recursive Spawning) within a running system. RSP instantiates the immutable parent chain as a runtime protocol, with a layered validation pipeline structurally identical to the BX3 Loop: the Purpose Layer validates intent, the Bounds Engine validates structural constraints, and the Fact Layer records the outcome. The three-layer architecture ensures that drift, orphaning, and unauthorized execution are not merely statistically unlikely outcomes but are structurally excluded from the set of architecturally possible states \cite{RecursiveSpawning_Sections_6_7_vFinal}.
